% !TEX program = xelatex
% 抽象代数视角下的初等数论讲义
% 编译方式：xelatex（需编译两遍以生成目录与交叉引用）
\documentclass[a4paper,11pt]{ctexart}

% ---- 数学与符号 ----
\usepackage{amsmath, amssymb, amsthm, mathtools}
\allowdisplaybreaks
\numberwithin{equation}{section}

% ---- 页面布局 ----
\usepackage[top=2.2cm,bottom=2.2cm,left=2.8cm,right=2.8cm]{geometry}
\usepackage{setspace}
\onehalfspacing

% ---- 页眉页脚 ----
\usepackage{fancyhdr}
\usepackage{lastpage}
\setlength{\headheight}{14.5pt}
\pagestyle{fancy}
\fancyhf{}
\fancyhead[L]{\leftmark}
\fancyhead[R]{\rightmark}
\fancyfoot[C]{\thepage~/~\pageref{LastPage}}
\fancypagestyle{plain}{
    \fancyhf{}
    \fancyfoot[C]{\thepage~/~\pageref{LastPage}}
    \renewcommand{\headrulewidth}{0pt}
}

% ---- 颜色与超链接 ----
\usepackage{xcolor}
\usepackage[colorlinks=true, allcolors=blue!45!black]{hyperref}

% ---- 列表 ----
\usepackage{enumitem}

% ---- 定理环境：编号跟随 section ----
\theoremstyle{plain}
\newtheorem{theorem}{Theorem}[section]
\newtheorem{lemma}[theorem]{Lemma}
\newtheorem{proposition}[theorem]{Proposition}
\newtheorem{corollary}[theorem]{Corollary}

\theoremstyle{definition}
\newtheorem{definition}[theorem]{Definition}
\newtheorem{example}[theorem]{Example}
\newtheorem{notation}[theorem]{Notation}

\theoremstyle{remark}
\newtheorem{remark}[theorem]{Remark}
\newtheorem{exercise}{习题}[section]

\renewcommand{\proofname}{Proof}

% ---- 自定义记号 ----
\newcommand{\Z}{\mathbb{Z}}
\newcommand{\Q}{\mathbb{Q}}
\newcommand{\R}{\mathbb{R}}
\newcommand{\C}{\mathbb{C}}
\newcommand{\F}{\mathbb{F}}
\newcommand{\N}{\mathbb{N}}
\DeclareMathOperator{\ord}{ord}
\DeclareMathOperator{\lcm}{lcm}
\DeclareMathOperator{\Img}{Im}
\newcommand{\leg}[2]{\left(\frac{#1}{#2}\right)}
\newcommand{\Ga}{\mathbb{Z}[i]}
\newcommand{\Gf}{\mathbb{Z}[\sqrt{-5}]}

\title{初等数论}
\author{}
\date{\today}

\begin{document}

\maketitle
\tableofcontents
\newpage

\setcounter{section}{-1}

\section{前言}

本文主要参考张贤科《初等数论》，并参考若干网络资料。

限于篇幅和主题，本文不从公理开始定义自然数、自然数的序和自然数的运算；下文直接使用通常的自然数、序、加法、乘法等概念。本文约定自然数包含 \(0\)。

\begin{notation}
    记全体整数为 \(\Z\)，全体自然数为 \(\N\)。
\end{notation}

下面不加证明地给出本文采用的自然数基本性质。

\begin{theorem}[数学归纳法]
    设 \(P\) 是 \(\N\) 上的一元谓词。若 \(P(0)\) 为真，且对任意自然数 \(n\)，\(P(n)\) 蕴含 \(P(n+1)\)，则对任意自然数 \(n\)，\(P(n)\) 为真。换言之，
    \[
    \bigl(P(0)\land \forall n(P(n)\to P(n+1))\bigr)
    \to \forall nP(n). 
    \]
\end{theorem}

\begin{corollary}[自然数良序]
    自然数的任何非空子集都有最小元。
\end{corollary}

由自然数的基本运算性质可得：

\begin{proposition}
    自然数关于加法构成交换幺半群，\(0\) 为加法幺元；关于乘法构成交换幺半群，\(1\) 为乘法幺元；乘法对加法满足分配律。因此 \((\N,+,\cdot,0,1)\) 构成交换半环。
\end{proposition}

网站 \url{https://adam.math.hhu.de/#/g/leanprover-community/nng4} 将这些命题的证明做成了分步游戏，感兴趣的读者可以自行探索。下文将不加证明地大量使用加法、乘法、幂运算的基本性质。

\begin{proposition}
    整数集 \(\Z\) 关于加法和乘法构成交换环 \((\Z,+,\cdot,0,1)\)，事实上是整环（没有非平凡零因子）。\(\Z\) 的乘法可逆元只有 \(\Z^\times=\{\pm1\}\)。
\end{proposition}

这里同样不给出证明。下文将进一步证明，\(\Z\) 是欧几里得整环，从而是主理想整环和唯一分解整环。

另外，由于存在保持加法和乘法的嵌入 \(\N \hookrightarrow \Z \hookrightarrow \Q \hookrightarrow \R \hookrightarrow \C\)，这里将结构在嵌入下的像直接作为结构本身，从而可以声称 \(\N \subseteq \Z \subseteq \Q \subseteq \R \subseteq \C\)。







\end{document}
